\section*{Введение}
\addcontentsline{toc}{section}{Введение}

Для обеспечения клиентам локальной сети доступа в Интернет необходимо настроить шлюз, который
принимает запросы клиентов и пересылает их во внешнюю сеть, а поступающие ответы
передаёт обратно во внутреннюю сеть. Такой шлюз реализуется с использованием механизма
\textbf{NAT} (Network Address Translation)~\cite{iptables_man}.

Целью данной практической работы является настройка виртуальной машины \texttt{gateway} на базе
 Ubuntu~20.04~LTS в качестве сетевого шлюза с двумя интерфейсами.

В рамках работы решаются следующие задачи:
\begin{itemize}
  \item настройка сетевых интерфейсов (netplan);
  \item включение передачи IP-пакетов (ip\_forward);
  \item настройка правил NAT через iptables (MASQUERADE);
  \item развёртывание и настройка DHCP-сервера \texttt{isc-dhcp-server}~\cite{isc_dhcp}.
\end{itemize}

Стенд: виртуальная машина \texttt{gateway} (VirtualBox),
интерфейсы \texttt{enp0s3}~(WAN/NAT) и \texttt{enp0s8}~(LAN, \texttt{192.168.\cfgLabStudentN.1/24});
клиентская виртуальная машина \texttt{Desktop1} (Ubuntu~Desktop).
